% !TeX root = ../../paper.tex

\subsection{Presence and Emotion in VR}
\label{sec:rw-presence-emotion}

Presence---the sense of being in the virtual place rather than merely looking at it---is the standard yardstick for experiences like ours. Slater's distinction between place illusion and plausibility~\cite{Slater.2009} maps well onto what a haptic proxy is supposed to deliver: touch that confirms vision should strengthen the plausibility of events, not just the sense of location. The most common way to measure presence remains the questionnaire. We use the Igroup Presence Questionnaire (IPQ)~\cite{Schubert.2001}, which separates general presence, spatial presence, involvement, and experienced realism. Questionnaires are a blunt instrument, and their sensitivity varies with the presence dimension being manipulated~\cite{presence-questionnaires}, but for a two-condition comparison of the same experience they remain the pragmatic choice.

A second strand of work examines how VR shapes emotional and physiological responses. Kobayashi et al.\ compared identically furnished real and virtual rooms and found that the VR versions produced higher arousal even though the real rooms were judged more pleasant~\cite{Kobayashi.2025}---VR is not emotionally neutral ground to begin with. Khundrakpam et al.\ used VR exercises to elicit distinct stress responses in a controlled way~\cite{Khundrakpam.2025}, Marocco et al.\ showed that arousal in VR is not simply negative, since physiologically activating nature scenes were experienced as restorative~\cite{Marocco.2025}, and Zhang et al.\ demonstrated that even the color of a virtual workspace shifts self-reported and physiological emotion~\cite{Zhang.2025}. For horror specifically, Lin studied a commercial zombie game and catalogued which elements of a VR horror experience players report as frightening and how they cope, including next-day effects~\cite{Lin.2017}.

Taken together, this work establishes that VR reliably provokes strong emotional responses and that presence and emotion are entangled. What it does not tell us is how a deliberately unpleasant tactile stimulus feeds into that loop. That gap is where this project sits: we hold the virtual content constant and vary only whether the player's hands find the expected horror or nothing at all.
